\usepackage{siunitx}
\usepackage{colortbl}
\usepackage{xspace}

\newcommand\accnum[1]{\num[round-precision=3]{#1}}
\newcommand\pernum[1]{\num[round-precision=1]{#1}}
\newcommand\scnum[1]{\num[scientific-notation=true,round-precision=1]{#1}}
\newcommand\numt[1]{\num[round-precision=2,round-mode=figures,scientific-notation=true]{#1}}
\newcommand\numo[1]{\num[round-precision=1,round-mode=figures,scientific-notation=true]{#1}}
%\newcommand{\numo}[1]{
%\pgfmathparse{#1}\pgfmathprintnumberto[precision=1, assume math mode=true]{\pgfmathresult}{\roundednumber}\roundednumber
%}
%\newcommand{\numt}[1]{
%\pgfmathparse{#1}\pgfmathprintnumberto[precision=2, assume math mode=true]{\pgfmathresult}{\roundednumber}\roundednumber
%}

\newcommand{\pref}[1]{page~\pageref{#1}}
\newcommand{\vect}[1]{\boldsymbol{#1}} % vector
\newcommand{\mat}[1]{\boldsymbol{#1}}  % matrix
\newcommand{\defeq}{\ensuremath{\mathrel{\mathop:}=}}
\newcommand{\eqdef}{\ensuremath{=\mathrel{\mathop:}}}

% \newcommand{\ipoint}[1]{\textit{\textbf{\color{darkgray}#1}}}
\newcommand{\ipoint}[1]{\textit{\textbf{\color{black!60}#1}}}
\newcommand{\todo}[1]{{\color{red}\textit{TODO}: #1}}
%\newcommand{\comment}[1]{{\color{blue}comment: #1}}
\newcommand{\commentout}[1]{}
\renewcommand{\b}[1]{{#1}}
\newcommand{\grbf}[1]{\mbox{\boldmath${#1}$\unboldmath}} %\renewcommand{\grbf}[1] {\mathbf{#1}}
\newcommand{\ns}[1]{\ensuremath{\mathbb{#1}}}
\newcommand{\fs}[1]{\ensuremath{\mathcal{#1}}}
\newcommand{\idiv}{\ensuremath{\nabla\cdot}}
\newcommand{\igrad}{\ensuremath{\nabla}}
\newcommand{\ilap}{\rotatebox[origin=c]{180}{$\nabla$}}
\newcommand{\icurl}{\ensuremath{\nabla \times}}
\newcommand{\tr}{\ensuremath{\operatorname{tr}}}
\newcommand{\id}{\ensuremath{\operatorname{id}}}
\newcommand{\ts}{\textsuperscript}
\newcommand{\half}[1]{\frac{#1}{2}}
\renewcommand{\d}[1]{\mathop{}\!\mathrm{d}#1}
\newcommand{\dt}{\d{t}}
\newcommand{\dx}{\d{\vect{x}}}
\newcommand{\op}[1]{\ensuremath{\mathcal{#1}}}
\newcommand{\di}[1]{\ensuremath{\mathbf{#1}}}


\newcommand{\MA}[1]{{\mathcal #1}}
\newcommand{\reals}{\mathbb{R}}
\newcommand{\codenm}[1]{{\tt  #1}}
\newcommand{\zapspace}{\topsep=0pt\partopsep=0pt\itemsep=0pt\parskip=0pt}

\newcommand{\msection}[1]{\section{#1}}
\newcommand{\msubsection}[1]{\subsection{#1}}
%\newcommand{\msection}[1]{\section{#1}}
%\newcommand{\msubsection}[1]{\subsection{#1}}


\definecolor{light-gray}{gray}{0.80}
%\renewcommand{\algorithmiccomment}[1]{#1}


\newcommand{\mcol}[2]{\multicolumn{#1}{r}{#2}}
\newcommand{\mrow}[2]{\multirow{#1}{*}{#2}}
\newcommand{\mrowrot}[2]{
\parbox[t]{2mm}{\multirow{#1}{*}{\rotatebox[origin=c]{90}{#2}}}
}


%\newcommand{\algcolor}[2]{\hspace*{-\fboxsep}\colorbox{#1}{\parbox{\linewidth}{#2}}}
\newcommand{\algcolor}[2]{\colorbox{#1}{\parbox{\linewidth}{#2}}}
\newcommand{\algemph}[1]{\algcolor{light-gray}{#1}}

\newcommand{\algcmt}[1]{\hfill {\footnotesize\ttfamily\textcolor{blue}{/* {#1} */}}}
%\newcommand{\algcmt}[1]{{\footnotesize\ttfamily\textcolor{blue}{\dotfill #1}}}
\newcommand{\algcc}[1]{\hfill {\footnotesize\ttfamily\textcolor{blue}{//{#1}}}}

% \let\oldparagraph\paragraph
% \renewcommand\paragraph{\subsubsection*}


\newcommand{\bfeval}[1]{
  \textcolor{black}{\textbf{\eval{#1}}}
}
\newcommand\aref{Algorithm~\ref}
\newcommand\appref{Appendix~\ref}
\newcommand\eref{Eq.~\ref}
\newcommand\fref{Fig.~\ref}
\newcommand\tref{Table~\ref}
\newcommand\sref{\S~\ref}
% \newcommand\sref{Section~\ref}
\newcommand\cmt[1]{\tcp*[r]{\scriptsize \color{gray!80!black}#1}}
\newcommand{\red}[1]{{\color{red} #1}}
\newcommand{\blue}[1]{{\color{blue} #1}}
\newcommand{\bluem}[1]{{\bf\color{blue} #1}}
\newcommand\ha{ \rowcolor{orange!0}}
\newcommand\hb{ \rowcolor{orange!15}}
\newcommand\hc{ \rowcolor{orange!40}}
\newcommand\hd{ \rowcolor{orange!45}}
\newcommand\he{ \rowcolor{orange!60}}
\newcommand\hf{ \rowcolor{orange!75}}
\newcommand\hg{ \rowcolor{orange!90}}

\newcommand\cha{\cellcolor{orange!0}}
\newcommand\chb{\cellcolor{orange!15}}
\newcommand\chc{\cellcolor{orange!40}}

\newcommand\ga{ \rowcolor{gray!0}}
\newcommand\gb{ \rowcolor{gray!15}}
\newcommand\gc{ \rowcolor{gray!40}}
\newcommand\gd{ \rowcolor{gray!45}}



% notation simplification
\newcommand\MS{{\mathcal{M}}}
\def\x{{\bf x}}
\def\H{{\bf H}}
\def\g{{\bf g}}
\def\Hg{{\bf Hg}}
% \def\0{{\bf 0}}
\def\s{{\bf s}}
\def\t{{\bf t}}
\def\p{{\bf p}}




\def\vv{{\bf v}}

\def\R{{\mathbb R}}

\def\AS{{\mathcal{A}}}
\def\V{{\mathbb{V}}}


\let\proof\relax
\let\endproof\relax
\usepackage{amsthm}
\let\lemma\relax
\let\assumption\relax
\let\remk\relax
\let\theorem\relax
\usepackage{amsthm}


\newtheorem{mytheorem}{Theorem}
\newtheorem{assumption}[mytheorem]{Assumption}
\newtheorem{lemma}[mytheorem]{Lemma}
\newtheorem{remk}[mytheorem]{Remark}
\newtheorem{definition}{Definition}

\newcommand{\es}[1]{{\color{red} #1}}

\newcommand{\MP}{\xspace\tiny{MP}}
\newcommand{\mixsolver}{\xspace{MixSolver}}
% \newcommand{\textapprox}{\raisebox{0.5ex}{\texttildelow}}



%\newtheorem{thm}{Theorem}
%\newtheorem{lem}[thm]{Lemma}
%\newtheorem{cor}[thm]{Corollary}
%\newtheorem{rem}[thm]{Remark}
%\newtheorem{remark}[thm]{Remark}
%\newtheorem{conj}[thm]{Conjecture}

\newcommand{\uniquepattern}{\emph{Unique}\xspace}
\newcommand{\replicatepattern}{\emph{Replicate}\xspace}
\newcommand{\shardpattern}{\emph{Shard}\xspace}
\newcommand{\shardpatternv}{\emph{Shard-V}\xspace}
\newcommand{\shardpatternh}{\emph{Shard-H}\xspace}
\newcommand{\partialpattern}{\emph{Partial}\xspace}
% \newcommand{\avgpattern}{\emph{params\_to\_average}\xspace}


\newcommand{\padfree}{\emph{StripPad}\xspace}
\newcommand{\union}{\emph{Union}\xspace}
\newcommand{\extract}{\emph{Extract}\xspace}
\newcommand{\genmeta}{\emph{UcpInfo}\xspace}
\newcommand{\load}{\emph{Load}\xspace}
\newcommand{\save}{\emph{Save}\xspace}

\newcommand{\pp}{pipeline parallelism\xspace}
\newcommand{\tp}{tensor parallelism\xspace}
\newcommand{\Pipe}{Pipeline}
\newcommand{\pipe}{pipeline}

\newcommand{\minjia}[1]{{\color{red}\small [Minjia: #1]}}
\newcommand{\kartik}[1]{{\color{blue}\small [Kartik: #1]}}
\newcommand{\ahan}[1]{{\color{orange}\small [Ahan: #1]}}
\newcommand{\jiahuan}[1]{{\color{orange}\small [Jiahuan: #1]}}
\newcommand{\kaidi}[1]{{\color{green}\small [Kaidi: #1]}}
\newcommand{\zihan}[1]{{\color{violet}\small [Zihan: #1]}}

\newcommand{\originalgrumbler}[2]{\begin{quote}\textcolor{blue}{\sl{\bf #1 says:} #2}\end{quote}}
\newcommand{\grumbler}[2]{\originalgrumbler{#1}{#2}}
\newcommand{\hide}[1]{}\newcommand{\draft}[1]{{\color{orange}#1}}
% \newcommand{\minjia}[1]{\grumbler{Minjia}{#1}}

% \newcommand{\name}{ElasticPD\xspace}
\newcommand{\name}{ADopt\xspace}
% \newcommand{\instance}{Instance-Level Elasticity\xspace}
% \newcommand{\cluster}{Cluster-Level Elasticity\xspace}
\newcommand{\instance}{FluidToken\xspace}
\newcommand{\cluster}{FluidRole\xspace}


\newcommand{\revisionA}[1]{{\color{blue} #1}}
\newcommand{\revisionB}[1]{{\color{red} #1}}
\newcommand{\revisionC}[1]{{\color{orange} #1}}
\newcommand{\revision}[1]{{\color{blue} #1}}

% ---- Revision marking for the design-section polish pass (Zihan) ----
% New text added in this pass (teal = unused by any author color):
\newcommand{\new}[1]{{\color{teal}#1}}
% Inline deletion of a short span (strikethrough; do NOT use across display math):
\newcommand{\del}[1]{{\color{gray}\sout{#1}}}
% Block deletion for paragraphs/equations: greys the content and fences it
% top and bottom with a full-width horizontal rule. Usage:
%   \begin{deleted} ...old text, equations ok... \end{deleted}
\newenvironment{deleted}
  {\par\noindent{\color{gray}\rule{\linewidth}{0.4pt}}\par\nobreak\color{gray}}
  {\par\noindent{\color{gray}\rule{\linewidth}{0.4pt}}\par}

\renewcommand{\emph}[1]{\textit{#1}}

%% Feedback macros.
\newcommand{\structure}[1]{{\color{blue} [structure: #1]}}
\newcommand{\important}[1]{{\color{red} [important: #1]}}

\newcommand{\wangjia}[1]{\textcolor{blue}{[Wangjia: #1]}}
\newcommand{\haochen}[1]{\textcolor{teal}{[Haochen: #1]}}
\newcommand{\ajay}[1]{\textcolor{brown}{[Ajay: #1]}}

\newcommand{\system}{\textsc{ADopt}\xspace}
\newcommand{\method}{\system}
\newcommand{\eg}{e.g.,\xspace}
\newcommand{\ie}{i.e.,\xspace}
\newcommand{\etal}{et al.\xspace}